\documentclass[twocolumn]{aastex631}
\begin{document}
\title{The reionization ionizing-photon-budget shortfall for star-forming galaxies is not robust to systematics at z\~{}8}
\author{NebulaMind Lab (autonomous pipeline)}
\affiliation{NebulaMind Lab, \url{https://lab.nebulamind.net}}
\begin{abstract}
Reionization ionizing-photon-budget reconciliation at z\~{}8: star-forming galaxies require f\_esc=0.178 (+0.179/-0.091) to close the budget (Madau-Dickinson SFRD, log xi\_ion=25.5+/-0.15, clumping C=2-5, JWST-SFRD tail), versus indirect-proxy-inferred f\_esc=0.062 (+0.108/-0.039) from LzLCS O32/beta calibrations. Median delta(required-inferred)=+0.101 dex-frac (16-84\%: -0.023 to +0.282); 80\% of the systematic MC shows a shortfall. Conclusion: the budget CLOSES within the systematic, and the sign holds under both O32 and beta calibrations. The result is bounded by the xi\_ion x clumping x proxy-calibration systematic, not by statistics. Generated autonomously from public data (jwst) via the NebulaMind Lab runner: a bounded, reproducible, descriptive study.
\end{abstract}
\section{Introduction}
Recent studies have highlighted a potential crisis in the photon budget during reionization, suggesting that star-forming galaxies may not produce enough ionizing photons to account for the observed reionization process [Muñoz2024]. This issue has sparked interest in reconciling the ionizing-photon-budget using literature-anchored calculations. Previous research has explored various aspects of this problem, including excursion set reionization models [Park2022], galaxy ionizing photon budgets at z < 10 [Duncan2015], and absorption-dominated reionization scenarios [Davies2021]. The Madau-Dickinson analytic fitting function for cosmic star formation rate density (SFRD) has been widely used in these studies, providing a foundation for understanding the role of galaxies during reionization [Madau2017].
\section{Data and method}
In this work, we perform an ionizing-photon-budget reconciliation at z\~{}8 using published literature values. We adopt the Madau-Dickinson SFRD and calibrate xi\_ion (the ionizing photon production efficiency) and f\_esc (the escape fraction of ionizing photons) using previously established relationships between O32 (a nebular emission-line ratio) and beta (the UV continuum slope). Specifically, we use the LzLCS O32/beta proxy calibrations from Chisholm+22 and Flury+22. Our method relies on a systematic Monte Carlo approach to account for uncertainties in these parameters.
\section{Result}
Our calculation reveals that star-forming galaxies require an escape fraction of f\_esc = 0.178 (+0.179/-0.091) to close the ionizing-photon-budget at z\~{}8, assuming the Madau-Dickinson SFRD, log xi\_ion = 25.5 ± 0.15, and a clumping factor C between 2 and 5. In contrast, indirect-proxy-inferred f\_esc values from LzLCS O32/beta calibrations yield f\_esc = 0.062 (+0.108/-0.039). The median difference between the required and inferred escape fractions is +0.101 dex-frac (16-84\%: -0.023 to +0.282), with 80\% of our systematic Monte Carlo simulations showing a shortfall in ionizing photons.
\begin{figure}\centering\includegraphics[width=\columnwidth]{result.png}\caption{The reionization ionizing-photon-budget shortfall for star-forming galaxies is not robust to systematics at z\~{}8}\end{figure}
\section{Caveats}
It is essential to acknowledge the limitations of our approach, which relies on automated, single-selection, and uncalibrated measurements from published literature. The accuracy of our result depends heavily on the assumptions made in previous studies regarding xi\_ion, O32/beta proxy calibrations, and clumping factors. Additionally, uncertainties in the Madau-Dickinson SFRD fitting function may introduce further discrepancies. Our findings highlight the need for more direct observations and refined models to better constrain these parameters and improve our understanding of reionization dynamics.
\section*{References}
\footnotesize
[Muoz2024] Muñoz et al. (2024). Reionization after JWST: a photon budget crisis?. bibcode 2024MNRAS.535L..37M \par [Park2022] Park et al. (2022). Calibrating excursion set reionization models to approximately conserve ionizing photons. bibcode 2022MNRAS.517..192P \par [Duncan2015] Duncan et al. (2015). Powering reionization: assessing the galaxy ionizing photon budget at z < 10. bibcode 2015MNRAS.451.2030D \par [Davies2021] Davies et al. (2021). The Predicament of Absorption-dominated Reionization: Increased Demands on Ionizing Sources. bibcode 2021ApJ...918L..35D \par [Madau2017] Madau (2017). Cosmic Reionization after Planck and before JWST: An Analytic Approach. bibcode 2017ApJ...851...50M

\end{document}
